\section{Related Works}
\label{sec:related-works}

Long-context efficiency has motivated a large body of work on efficient attention, which can be broadly grouped into two directions: replacing dense softmax attention with cheaper linear or recurrent alternatives, and retaining softmax attention while restricting its receptive field.
Linear attention~\citep{katharopoulos2020linear,choromanski2020performer} replaces the softmax kernel with a linear-complexity surrogate, while state-space models such as Mamba~\citep{gu2023mamba} replace attention with a selective recurrence over hidden states.
Hybrid stacks~\citep{minimax01,minimaxm1} interleave linear blocks with full-attention blocks, reducing the number of quadratic layers while preserving part of the exact-softmax capacity.
Fixed-pattern attention keeps softmax attention but imposes a predefined support, including local windows, global tokens~\citep{beltagy2020longformer,zaheer2020bigbird}, and attention sinks with a sliding window~\citep{xiao2023streaming}.
These approaches reduce long-context cost either by replacing dense attention in part or in full, or by using a content-agnostic attention pattern.

Beyond fixed sparse patterns, adaptive sparse attention makes the attended support depend on the input. 
Existing methods differ mainly in when this support is constructed and whether the selector is trained as part of the model. 
Inference-time sparsification operates on a pretrained Full-Attention backbone and constructs sparse supports only during serving. H2O~\citep{zhang2023h2o} and SnapKV~\citep{li2024snapkv} prune the KV cache during decoding using accumulated attention statistics, Quest~\citep{tang2024quest} performs page-level importance estimation per query, MInference~\citep{jiang2024minference} and FlexPrefill~\citep{lai2025flexprefill} dispatch per-head sparse kernels at prefill, and InfLLM~\citep{xiao2024infllm} maintains attention sinks, a local context window, and retrievable chunks. These methods inherit the training cost of Full Attention and leave at least one inference phase near Full-Attention speed. Natively trained sparse-attention designs train the indexer during pretraining and form the closest prior work to \method{}. NSA~\citep{yuan2025nsa} targets MQA/MHA backbones with three parallel branches: compressed attention for coarse global context, selected attention over fine-grained blocks, and a sliding window for local context. InfLLM-V2~\citep{zhao2025infllmv2} achieves zero-shot dense-to-sparse switching by unifying parameter-free block selection with a local sliding window. MoBA~\citep{lu2025mobamixtureblockattention} also operates on GQA but uses very large KV blocks scored by block-averaged keys, and trains its indexer only through the language-modeling gradient. DSA~\citep{deepseekai2025deepseekv32pushingfrontieropen} sits on top of MLA in its MQA mode: a multi-head ReLU-based lightning indexer scores tokens individually, all query heads share a single Top-$k$ index, and selection is token-level.
\method{} differs from this neighborhood along two axes that are taken up together: per-GQA-group Top-$k$ sharing combined with block-level selection, which gives multi-group block-granular retrieval while keeping KV reads contiguous.

Efficient kernels are essential for sparse attention to translate theoretical FLOP reduction into wall-clock speedups. FlashAttention~\citep{dao2022flashattention} and FlashAttention-2~\citep{dao2023flashattention2} introduced IO-aware tiled softmax attention, and FlashDecoding~\citep{dao2023flashdecoding} extended this to memory-bound decoding. Open-source block-sparse kernels such as Flash-Sparse-Attention~\citep{yan2025fsa} and FlashMoBA~\citep{xiao2025flashmoba} have made block-sparse variants of this recurrence available. \method{}'s kernel, described in \cref{sec:infrastructure}, reuses the FlashAttention algorithmic skeleton with a loop ordering tuned to the GQA-native, block-granular access pattern \method{} produces.
